\documentclass[a4paper,11pt]{article}

\newcommand{\TPnumero}{Lab 5 -- Solution}
% =========================================================
% Tutorial / lab preamble — Time Series Analysis module
% article class + fancyhdr + tcolorbox + listings (English)
% Author : Aymen Ben Brik (aymen.benbrik@esprit.tn)
% =========================================================

\usepackage[utf8]{inputenc}
\usepackage[T1]{fontenc}
\usepackage[english]{babel}
\usepackage{lmodern}
\usepackage{amsmath, amssymb, amsthm, mathtools, bm}
\usepackage{graphicx}
\usepackage{booktabs}
\usepackage{array}
\usepackage{multirow}
\usepackage{colortbl}
\usepackage{xcolor}
\usepackage{tikz}
\usetikzlibrary{positioning, arrows.meta, calc, shapes, fit, backgrounds, matrix}
\usepackage{listings}
\usepackage[most]{tcolorbox}
\usepackage{geometry}
\geometry{a4paper, margin=2cm, headheight=15pt}
\usepackage{fancyhdr}
\usepackage[hidelinks]{hyperref}
\usepackage{enumitem}

% --- Colours ---
\definecolor{espritBlue}{RGB}{30, 60, 110}
\definecolor{espritAccent}{RGB}{200, 80, 30}
\definecolor{boxEx}{RGB}{30, 60, 110}
\definecolor{boxQ}{RGB}{180, 120, 0}
\definecolor{boxC}{RGB}{30, 100, 60}
\definecolor{boxAtt}{RGB}{200, 80, 30}

% --- Listings ---
\definecolor{codebg}{RGB}{248,248,248}
\definecolor{codeframe}{RGB}{220,220,220}
\definecolor{keyword}{RGB}{0,82,155}
\definecolor{comment}{RGB}{88,110,117}
\definecolor{string}{RGB}{133,153,0}

\lstdefinestyle{pythonstyle}{
  language=Python,
  basicstyle=\ttfamily\small,
  keywordstyle=\color{keyword}\bfseries,
  commentstyle=\color{comment}\itshape,
  stringstyle=\color{string},
  backgroundcolor=\color{codebg},
  frame=single,
  rulecolor=\color{codeframe},
  frameround=tttt,
  breaklines=true,
  showstringspaces=false,
  tabsize=2
}
\lstdefinestyle{Rstyle}{
  language=R,
  basicstyle=\ttfamily\small,
  keywordstyle=\color{keyword}\bfseries,
  commentstyle=\color{comment}\itshape,
  stringstyle=\color{string},
  backgroundcolor=\color{codebg},
  frame=single,
  rulecolor=\color{codeframe},
  frameround=tttt,
  breaklines=true,
  showstringspaces=false,
  tabsize=2
}
\lstset{style=pythonstyle}

% --- Common macros (configurable path) ---
\providecommand{\macrosPath}{../../preamble/macros.tex}
\input{\macrosPath}

% --- Exercise counter ---
\newcounter{excounter}
\renewcommand{\theexcounter}{\arabic{excounter}}

\newtcolorbox[use counter=excounter]{exercise}[1][]{
  enhanced, breakable,
  colback=boxEx!5, colframe=boxEx, fonttitle=\bfseries,
  title={Exercise~\theexcounter\ifstrempty{#1}{}{ -- #1}},
  arc=2pt, boxrule=0.6pt
}
\newtcolorbox{question}[1][]{
  enhanced, breakable,
  colback=boxQ!5, colframe=boxQ, fonttitle=\bfseries,
  title={Question\ifstrempty{#1}{}{ -- #1}},
  arc=2pt, boxrule=0.5pt
}
\newtcolorbox{solution}[1][]{
  enhanced, breakable,
  colback=boxC!5, colframe=boxC, fonttitle=\bfseries,
  title={Solution\ifstrempty{#1}{}{ -- #1}},
  arc=2pt, boxrule=0.5pt
}
\newtcolorbox{warning}[1][]{
  enhanced, breakable,
  colback=boxAtt!5, colframe=boxAtt, fonttitle=\bfseries,
  title={Warning\ifstrempty{#1}{}{ -- #1}},
  arc=2pt, boxrule=0.5pt
}

% --- Headers / footers ---
\pagestyle{fancy}
\fancyhf{}
\fancyhead[L]{\textsc{\moduleNom} -- \auteurNom}
\fancyhead[R]{\TPnumero}
\fancyfoot[C]{\thepage}
\fancyfoot[L]{\auteurInstitution}
\fancyfoot[R]{\auteurEmail}
\renewcommand{\headrulewidth}{0.4pt}
\renewcommand{\footrulewidth}{0.2pt}

% --- Placeholder ---
\providecommand{\TPnumero}{Lab}


\title{Lab 5: ARMA estimation, diagnostics and forecasting\\
\large \emph{Reference solution}}
\author{\auteurNom\\\auteurEmail\\\auteurInstitution}
\date{\anneeUniversitaire}

\begin{document}
\maketitle

\noindent\emph{Reference Python code in \texttt{correction.py}.
All numerical values below are obtained with}
\texttt{numpy.random.seed(42)}.

\setcounter{excounter}{0}

% =========================================================
\begin{exercise}[AR$(2)$ simulation and estimation]\end{exercise}

\begin{solution}
The PACF of the simulated series cuts off after lag $2$, suggesting
AR$(2)$. Fit comparison:
\begin{center}
\begin{tabular}{lrrrr}
\toprule
order        & log-lik. & AIC      & BIC      & LB$(24)$ $p$ \\
\midrule
$(1,0,0)$    & $-720.28$ & $1446.55$ & $1459.19$ & $0.0000$ \\
$(2,0,0)$    & $-684.91$ & $\bm{1377.81}$ & $\bm{1394.67}$ & $0.7506$ \\
$(3,0,0)$    & $-684.84$ & $1379.68$ & $1400.76$ & $0.7498$ \\
\bottomrule
\end{tabular}
\end{center}
The MLE on $(2, 0, 0)$ gives
$\widehat\phi_1 = +0.7051$, $\widehat\phi_2 = -0.3645$, close to the
true $(0.6, -0.3)$ given $T = 500$. AR$(1)$ is rejected by both AIC
and Ljung--Box; AR$(3)$ has a marginally larger BIC --- AR$(2)$ is
the parsimonious choice.
\end{solution}

% =========================================================
\begin{exercise}[MA$(1)$ identification]\end{exercise}

\begin{solution}
Empirical $\widehat\rho(h)$ for $h = 1, \dots, 5$:
\[
  (0.517,\ 0.024,\ -0.059,\ -0.073,\ -0.069),
\]
the ACF cuts off at lag~$1$. The theoretical $\rho(1) = 0.7 /
1.49 = 0.4698$ is recovered to within $\pm 0.05$.

Fit comparison:
\begin{center}
\begin{tabular}{lrrrr}
\toprule
order        & log-lik. & AIC      & BIC      & LB$(24)$ $p$ \\
\midrule
$(0,0,1)$    & $-688.01$ & $\bm{1382.02}$ & $\bm{1394.67}$ & $0.6726$ \\
$(1,0,0)$    & $-732.61$ & $1471.22$ & $1483.86$ & $0.0000$ \\
\bottomrule
\end{tabular}
\end{center}
$\widehat\theta = +0.7558$ (true $0.7$). The AR$(1)$ alternative
fails Ljung--Box because an MA$(1)$ has only one non-zero
autocorrelation, while AR$(\infty)$ is needed to reproduce that
single spike via an autoregressive model: an AR$(1)$ approximation
leaves persistence on every lag, hence the residual autocorrelation.
\end{solution}

% =========================================================
\begin{exercise}[ARMA$(1, 1)$ on simulated data]\end{exercise}

\begin{solution}
\begin{center}
\begin{tabular}{lrrrr}
\toprule
order        & log-lik. & AIC      & BIC      & LB$(24)$ $p$ \\
\midrule
$(1,0,0)$    & $-720.66$ & $1447.31$ & $1459.96$ & $0.0000$ \\
$(0,0,1)$    & $-728.47$ & $1462.93$ & $1475.58$ & $0.0000$ \\
$(1,0,1)$    & $-687.45$ & $1382.91$ & $\bm{1399.77}$ & $0.6568$ \\
$(2,0,2)$    & $-684.82$ & $\bm{1381.64}$ & $1406.93$ & $0.8687$ \\
\bottomrule
\end{tabular}
\end{center}
AIC marginally favours $(2, 0, 2)$, BIC favours $(1, 0, 1)$. The
$(2, 0, 2)$ model adds three extra parameters for a $1.27$-point
gain in log-likelihood --- classic over-fitting (AR/MA cancellation
between the spurious factors). The parsimonious choice is
$(1, 0, 1)$ with $\widehat\phi = +0.5093$ and
$\widehat\theta = +0.4730$, very close to the true $(0.5, 0.4)$.
\end{solution}

% =========================================================
\begin{exercise}[Box--Jenkins on the Atlanta residual]\end{exercise}

\begin{solution}
On the Atlanta decomposition residual ($n = 1644$, sd $= 2.884$):
\begin{center}
\begin{tabular}{lrrrrr}
\toprule
order        & log-lik. & AIC      & BIC      & $\widehat\sigma^2$ & LB$(24)$ $p$ \\
\midrule
$(1,0,0)$    & $-4071.62$ & $8149.24$ & $8165.46$ & $8.293$ & $0.0000$ \\
$(2,0,0)$    & $-4065.35$ & $8138.70$ & $8160.32$ & $8.230$ & $0.0000$ \\
$(0,0,1)$    & $-4071.22$ & $8148.44$ & $8164.65$ & $8.289$ & $0.0000$ \\
$(1,0,1)$    & $-3928.70$ & $7865.39$ & $7887.01$ & $6.945$ & $0.0000$ \\
$(2,0,1)$    & $-3871.89$ & $\bm{7753.78}$ & $\bm{7780.81}$ & $6.480$ & $0.0000$ \\
\bottomrule
\end{tabular}
\end{center}
ARMA$(2, 1)$ wins both AIC and BIC. The Ljung--Box test still
rejects, however: the residual carries structure beyond what a
linear ARMA can capture (residual seasonality / volatility
clusters --- to be addressed in Chapter~6).

Best parsimonious choice: \textbf{ARMA$(2, 1)$} as a useful
approximation, with the caveat that the model is not a perfect
fit.
\end{solution}

% =========================================================
\begin{exercise}[Forecasting]\end{exercise}

\begin{solution}
With ARMA$(2, 1)$, the $24$-step forecast (figure
\texttt{lab5\_ex5\_forecast.png}) returns smoothly to the
unconditional mean of zero, and the $95\%$ confidence band widens
toward
\[
  \pm 1.96\, \widehat\sigma\,
        \sqrt{\sum_{j = 0}^{\infty} \psi_j^2}
  \;=\; \pm 1.96\, \cdot \,\text{(unconditional sd of $X_t$)}.
\]

\textbf{Hold-out validation.} Splitting at $80\%$ ($n_{\text{train}}
= 1315$, $n_{\text{test}} = 329$): out-of-sample RMSE $= 2.699$,
to be compared with the in-sample residual sd $= 2.598$. The
out-of-sample error is only $\approx 4\%$ larger than the in-sample
sd: the model generalises reasonably well, but the seasonal /
volatility structure noted in Exercise~$4$ shows up as small
systematic forecast errors.
\end{solution}

\end{document}
